\section{Local periods and an effective inverse neighbourhood}
\label{sec:local-periods}

We construct local period coordinates with quantitative evaluation data.
The central coefficients and the six transverse coefficient directions
are algebraic. The inverse period section will generally have
transcendental coefficients. We therefore distinguish exact algebraic
operations from the evaluation of analytic functions throughout the
construction.

\subsection{Evaluation data}

A name for a complex number supplies rational complex approximations
with any prescribed positive absolute error. Statements about an
arbitrary named parameter are relative to this access; they do not
assert that every complex number is computable. For computable parameters
the procedures below are ordinary algorithms.

\begin{definition}\label{def:bounded-analytic}
A bounded analytic presentation of a holomorphic function consists of
a polydisc with specified positive rational radii, an evaluation procedure,
and a known finite upper bound on a slightly larger closed polydisc.
All derivative bounds are taken on specified strictly smaller domains.
\end{definition}

The larger complex domain is essential. Cauchy's formula both bounds
derivatives and evaluates them by contour quadrature. We do not infer
derivative access from an arbitrary pointwise evaluator.

Let $c_*$ be an algebraic coefficient triple defining the eight-node
central complete intersection in $\PP^3\times\PP^3$. Fix six algebraic
coefficient directions for which the local period derivative has rank
six, and write $s\in\CC^6$ for their coordinates. The resulting ordered
triple of equations $F(s)$ has bidegrees $(1,1),(1,1),(2,2)$.
Let $\Psi$ be its ordered relative residue form. The node labels and
oriented block relations are those of Section~\ref{sec:algebraic-family};
retain its space $R$ in \eqref{eq:period-space}.
The map $\pi_R:R\longrightarrow\CC^6$ retains the first three
coordinates in each block. Its inverse fills the fourth coordinate
with the negative sum of the other three.

\begin{theorem}\label{thm:evaluated-periods}
From these algebraic data one can construct a positive rational radius
$\rho$, a ball $B_{\rm ev}=\{p\in R:\|p\|<\rho\}$, and a
holomorphic coefficient map $\sigma$ with the following properties.
The fibres
\[
 X_p=Z\bigl(F(\sigma(\pi_Rp))\bigr)
\]
have local periods exactly $p_i$ for the form $\Psi_p$. Their
discriminant is exactly $\bigcup_i\{p_i=0\}$, and every singularity
of a fibre is one of the prescribed ordinary double points.
The family is flat and projective, its total space is smooth, its
polarization is $\OO(2,2)$, and its relative dualizing form is $\Psi$.
The coefficient functions, local normal forms, residue units and
derivatives of every prescribed finite order have bounded analytic
presentations on specified smaller domains.

One can also select an algebraic coefficient point $s_*$ whose fibre
is smooth and whose evaluated period point $p_*$ belongs to
$B_{\rm ev}^{\circ}=B_{\rm ev}\setminus\bigcup_i\{p_i=0\}$,
with $s_* = \sigma(\pi_Rp_*)$.
\end{theorem}

No comparison between the nonzero coordinates $p_i$ occurs in this
statement. The theorem supplies a known observation domain; it does
not give a complexity bound as a period approaches zero.

\subsection{An implicit function with bounds}

We first give the analytic operation used in the proof. Suppose that
$G(y,z)$ has a bounded analytic presentation, $G(0,0)=0$, and
$A=D_yG(0,0)$ is invertible, with an evaluated inverse and a known
upper bound for its norm. Use the maximum norm and its induced
operator norm. Bounds on the second derivatives allow us to find
positive rational $r,\eta$ such that, for $\|y\|\le r$ and
$\|z\|\le\eta$,
\begin{equation}\label{eq:implicit-bounds}
 \|I-A^{-1}D_yG(y,z)\|\le q=\tfrac14,
 \qquad \|A^{-1}G(0,z)\|\le r/4.
\end{equation}
One may first require strict stronger bounds, so their verification
does not require an equality test on a named complex number.
The map $T_z(y)=y-A^{-1}G(y,z)$ sends the closed $r$-ball into
the $r/2$-ball and is a $q$-contraction. Its unique fixed point satisfies
\begin{equation}\label{eq:implicit-tail}
 \|y(z)-T_z^n(0)\|
       \le \frac{q^n}{1-q}\,\frac r4.
\end{equation}
Choose $n$ to control this tail, then evaluate the finitely many
compositions with errors whose propagated sum is smaller than the
remaining tolerance. The interior margin keeps these compositions
in their specified domains.

Computable local inversion with derivative access is classical; see
\cite[Theorem~13 in the arXiv-v2 text]{Ziegler}. Here the displayed
contraction estimates provide the additional complex-domain bounds
used by the subsequent period construction.

Uniform convergence of the holomorphic iterates proves holomorphy of
$y(z)$ and gives the bound $\|y\|\le r/2$. Cauchy's formula on
an interior product contour evaluates any prescribed derivative.
A bound on a larger interior contour controls the variation of the
integrand, and therefore the error of its Riemann sums. This proves
both derivative evaluation and derivative bounds. Sums, products,
compositions with a known image margin, and reciprocals with a known
positive denominator bound preserve these properties.

\subsection{Parameter-preserving holomorphic Morse coordinates}

At an algebraic central node, choose a nonzero two-equation Jacobian
minor. The preceding construction eliminates two affine variables,
leaving a bounded holomorphic equation $f(z,s)=0$ in four variables
$z$. Its central Hessian is nonsingular. A second implicit solution,
now of $D_zf=0$, gives the critical point $c(s)$. Set
$\tau(s)=-f(c(s),s)$. Taylor's formula gives
\begin{equation}\label{eq:integrated-hessian}
 \begin{split}
 f(c(s)+v,s)-f(c(s),s)&=v^TH(v,s)v,\\
 H(v,s)&=\int_0^1(1-t)D_z^2f(c(s)+tv,s)\,\dd t.
 \end{split}
\end{equation}
The integral is evaluated on a fixed real interval using a bound on
the next derivative. Thus $H$ is a bounded evaluated holomorphic
symmetric matrix. Its central value is algebraic and nonsingular.
Symmetric elimination over the algebraic numbers gives $R_0$ with
$R_0^TH(0,0)R_0=I$.

Writing $v=R_0u$, put
$K(u,s)=R_0^TH(R_0u,s)R_0-I$. On a known smaller domain
$\|K\|\le1/4$. The series
\begin{equation}\label{eq:morse-square-root}
 C(u,s)=\sum_{j\ge0}\binom{1/2}{j}K(u,s)^j
\end{equation}
has remainder at most $q^{N+1}/(1-q)$ after degree $N$, where
$q=1/4$. Each power of the symmetric matrix $K$ is symmetric,
and the convergent series identity gives $C^TC=C^2=I+K$.
Hence $u^T(I+K)u=(Cu)^T(Cu)$. The map $u\mapsto C(u,s)u$
has derivative $I$ at the central point. Applying
\eqref{eq:implicit-bounds} once more supplies an inverse with the
parameter $s$ unchanged. In the resulting coordinates the equation is
\begin{equation}\label{eq:quadric-normal-form}
             Q(w)=w_1^2+w_2^2+w_3^2+w_4^2=\tau(s).
\end{equation}
All maps and inverses have known domains and derivative bounds.
For the convergent analytic splitting lemma without these evaluation
data, see \cite[Theorem~2.4 in the arXiv-v2 text]{GreuelPfister}.

The ordered complete-intersection residue becomes
\begin{equation}\label{eq:local-residue}
 \Psi_s=a(w,s)\operatorname{Res}
   \frac{\dd w_1\wedge\cdots\wedge\dd w_4}{Q(w)-\tau(s)}.
\end{equation}
The first elimination divides by the chosen equation minor; subsequent
changes multiply by evaluated Jacobian determinants. The minor is
bounded away from zero on a smaller domain. Thus $a$ has a bounded
analytic presentation and $a(0,0)$ is nonzero and algebraic. The
fixed local orientation can be included in the sign of $a$.

\subsection{A period series valid at the discriminant}

Choose a known rational $r>0$ and a parameter polydisc such that
$a$ is holomorphic near their product with $\|w\|_\infty\le r$,
and $|a|\le M$. Let $\alpha_0$ be half the standard oriented
volume form on the unit sphere $S^3$, so $\int_{S^3}\alpha_0=\pi^2$.
The function
\[
 A(q,s)=\int_{S^3}a(qy,s)\alpha_0(y)
\]
is holomorphic and even in $q$: the antipodal map of $S^3$ has
degree $+1$. Consequently
\begin{equation}\label{eq:even-period-series}
 A(q,s)=\sum_{k\ge0}b_{2k}(s)q^{2k},
 \qquad |b_{2k}(s)|\le\pi^2Mr^{-2k}.
\end{equation}
The bound is Cauchy's estimate on the $q$-circle.

These coefficients can be evaluated without integrating a shrinking
sphere. Write $a(w,s)=\sum_\alpha a_\alpha(s)w^\alpha$.
Odd sphere moments vanish by reflection. Separating four Gaussian
integrals and then using polar coordinates gives, for $|m|=k$,
\[
 \int_{S^3}y^{2m}\alpha_0
 =\frac{\prod_{j=1}^4\Gamma(m_j+1/2)}{\Gamma(k+2)}
 =\frac{\pi^2}{(k+1)!}
       \prod_{j=1}^4\frac{(2m_j)!}{4^{m_j}m_j!}.
\]
It follows that
\begin{equation}\label{eq:period-moments}
 b_{2k}(s)=\frac{\pi^2}{(k+1)!}
   \sum_{|m|=k}a_{2m}(s)
           \prod_{j=1}^4\frac{(2m_j)!}{4^{m_j}m_j!}.
\end{equation}
Each coefficient $a_{2m}$ is a derivative divided by its factorial
and is evaluated by the Cauchy procedure above.

Define $U(t,s)=\sum_{k\ge0}b_{2k}(s)t^k$. For $|t|\le r^2/4$,
\begin{equation}\label{eq:unit-series-tail}
 \left|U(t,s)-\sum_{k=0}^Nb_{2k}(s)t^k\right|
 \le\pi^2M\frac{(|t|/r^2)^{N+1}}{1-|t|/r^2}
 \le\frac43\pi^2M\,4^{-(N+1)}.
\end{equation}
For $\tau\ne0$, parametrizing the local sphere by $w=qy$,
$q^2=\tau$, multiplies the standard residue by $q^2$.
The local period therefore has the holomorphic continuation
\begin{equation}\label{eq:period-unit}
                       P(s)=\tau(s)U(\tau(s),s).
\end{equation}
Either square root gives the same integral. Formula
\eqref{eq:period-unit} evaluates the period even at $\tau=0$;
it neither tests whether a named number is zero nor divides by $\tau$.

Since $U(0,0)=\pi^2a(0,0)\ne0$, evaluation of its nonzero
central value and derivative bounds give a further known shrink on
which
\begin{equation}\label{eq:known-unit-bounds}
          0<c\le |U(\tau(s),s)|\le C.
\end{equation}
Repeating the construction at eight nodes and taking finitely many
minima supplies one common coefficient neighbourhood, with the same
evaluation properties for all the periods and their derivatives.

\subsection{Excluding other singularities}

Choose disjoint coordinate boxes about the eight algebraic nodes,
with algebraic centres and rational side lengths. Their closures lie
inside the elimination and inverse Morse domains. Shrink the boxes
so that every common zero in each closure belongs to the uniquely
eliminated branch; the strict implicit-function margins preserve this
property under a known coefficient change.

On the complement of these boxes in the compact ambient product use
the fixed Fubini--Study metrics and connections. If
$A_0=\nabla_xF(0)$, then
\[
        g_0(x)=|F(0,x)|^2+\lambda_{\min}(A_0A_0^*)
\]
is strictly positive there. A zero would be a singular common zero
outside the eight boxes. A positive rational lower bound $m$ is found by
successive finite meshes and known variation bounds. The central
compact domains and coefficients are algebraic, so their coverage
conditions can also be checked by real algebraic quantifier elimination.
Polynomial dependence on $s$ and the eigenvalue perturbation inequality
give a known coefficient radius on which $g_s\ge m/2$. Thus no new
singularities occur on that complement. In each node box,
\eqref{eq:quadric-normal-form} shows that a singularity occurs
precisely when $\tau_i=0$ and is then an ordinary double point.

At every fibre point the three equations form a regular sequence.
At a smooth point this follows from their independent differentials;
at a node two equations eliminate variables and the third is the
nonzero quadratic hypersurface equation. The fixed-degree Koszul
resolution gives a constant Hilbert polynomial and flatness over the
coefficient neighbourhood. Adjunction gives the relative form $\Psi$.
The same Koszul resolution gives
$H^0(\OO_{X_s})=\CC$ and $H^1(\OO_{X_s})=H^2(\OO_{X_s})=0$.
The fibres are normal, since they are local complete intersections
with at most isolated nodes. Connected normal fibres are integral.
The line bundle $\OO(2,2)$ remains relatively ample.

\subsection{Period inversion and an algebraic reference fibre}

The two sphere relations imply $P(s)\in R$. Set
$Q_0(s)=\pi_RP(s)$. Its central derivative is the prescribed
invertible matrix, with entries in $\pi^2\overline{\QQ}$.
Apply the bounded implicit construction to $Q_0(s)-v=0$.
It gives a holomorphic $s=\sigma(v)$ on a known neighbourhood, with
\begin{equation}\label{eq:inverse-period-identity}
                    P(\sigma(\pi_Rp))=p.
\end{equation}
Choose a rational ball $B_{\rm ev}\subset R$ whose image under
$\pi_R$ lies strictly inside that neighbourhood and whose image under
$\sigma\pi_R$ lies strictly inside the geometric coefficient domain.
The computed image bounds give this choice. Equations
\eqref{eq:period-unit}--\eqref{eq:known-unit-bounds} show that
\[
 \tau_i(\sigma(\pi_Rp))=p_i/U_i(p),\qquad
                  0<c\le|U_i(p)|\le C.
\]
This proves the exact discriminant assertion. At a nodal point
$\dd\tau_i=U_i^{-1}\dd p_i$; the coordinate functional $p_i$
is nonzero on $R$. The total space is therefore smooth there,
and it is smooth over every regular fibre point as well.

There is also a known small coefficient ball on which
$\sigma(Q_0(s))=s$ and $P(s)\in B_{\rm ev}$: shrink until the
derivative bound for $Q_0$ puts its image inside the inverse domain,
and use uniqueness of the implicit solution. Enumerate algebraic
points in this ball. The smooth locus is nonempty and open, because
the period map is locally invertible and no coordinate $p_i$ vanishes
identically on $R$. Algebraic points are dense. At an algebraic
candidate, smoothness is decided by the projective polynomial
Jacobian rank tests. The search therefore finds a smooth $s_*$.
The period $p_*=P(s_*)$ is evaluated by the series above and satisfies
the assertions of Theorem~\ref{thm:evaluated-periods}.

The reference fibre has algebraic equations, although its period point
need not be algebraic. This is the fibre on which we use exact
semialgebraic topology. At general period points we use analytic
evaluations with strict margins, without applying algebraic sign
tests to transcendental coefficients.
